%% ma: L10-B02-0029
%% lop: 10
%% chuong: 1
%% bai: 2
%% dang: 10.2.3
%% loai: TN4
%% muc: TH
%% dap_an: "C"
%% nguon: "Ảnh giáo viên gửi 10/10/2026 (ThS. Nguyễn Hoàng Việt, Chương 2 Tập hợp), Đề 2 (đề thứ hai, không ghi tiêu đề), Phần 1 Câu 6"
%% kien_thuc: "Bài 2 (tập hợp và các phép toán trên tập hợp)"
%% trang_thai: chinh-thuc
\begin{ex}
Trong các tập hợp sau, tập nào là tập rỗng?
\choice{$\{x\in\mathbb R\mid x^2+5x-6=0\}$}{$\{x\in\mathbb Q\mid3x^2-5x+2=0\}$}{\True $\{x\in\mathbb Z\mid x^2+x-1=0\}$}{$\{x\in\mathbb R\mid x^2+5x-1=0\}$}
\loigiai{A: $x^2+5x-6=(x+6)(x-1)=0$ có nghiệm thực. B: $3x^2-5x+2=(3x-2)(x-1)=0$ có nghiệm hữu tỉ $\dfrac23,1$. C: $x^2+x-1=0$ có nghiệm $\dfrac{-1\pm\sqrt5}{2}$ vô tỉ, không có nghiệm nguyên nên tập rỗng. D: $\Delta=29>0$, có nghiệm thực.}
\end{ex}
